\problem{CA01}{Weighted Sampling from a Stream}{}
Items arrive sequentially with positive weights $w_i$. The total number of items is unknown. Only one candidate item and a constant number of numerical values may be stored. Design an algorithm such that after any $t\ge1$ arrivals, item $i$ is the candidate with probability $w_i/\sum_{j=1}^t w_j$. Each update may use an independent continuous uniform random variable. Prove the claim and explain why replacing the candidate whenever the new item is heavier fails.
\solution Maintain the cumulative weight $W$. Select the first item and set $W=w_1$. When item $t$ arrives, first update $W\leftarrow W+w_t$, then replace the candidate with probability $w_t/W$. Otherwise retain it. Each update takes constant arithmetic time and stores only the candidate identifier and cumulative weight.

For the inductive proof, let $W_{t-1}$ denote the previous total. An old item $i<t$ survives precisely when it was selected previously and the new update does not replace it. Independence of the fresh random variable gives
\[
\Prob(\text{candidate}=i)
=\frac{w_i}{W_{t-1}}\left(1-\frac{w_t}{W_t}\right)
=\frac{w_i}{W_t}.
\]
The new item is selected with probability $w_t/W_t$ directly. The base case $t=1$ is immediate, proving the invariant at every fixed length. It also applies when the ending length is fixed externally or independent of the sampling randomness.

For weights $1,2,3$, the second update selects item 2 with probability $2/3$ and the third replaces with probability $1/2$. Final probabilities are $1/6,1/3,1/2$. Deterministic replacement by a heavier item always retains the last item in a strictly increasing stream, so it cannot have the required distribution.

The proof applies to a fixed sequence of weights. Random weights independent of the sampling randomness can be handled by conditioning on the entire sequence. If future weights adapt to publicly revealed candidate choices, that conditioning argument is no longer valid without further analysis. The space bound counts stored values; their exact bit lengths may grow.
\insight The candidate represents all past weight through a probability invariant. A new item competes against the entire accumulated weight, not merely against the weight of the current representative.

\problem{CA02}{The Median of Two Sorted Arrays}{}
Two nondecreasing real arrays $A,B$ have lengths $m,n$ with $m+n>0$. Random access is allowed, but fully merging the arrays is not. Find their combined median, averaging the central two elements when the total length is even. Give an $O(\log(\min\{m,n\}+2))$-time, $O(1)$-extra-space algorithm covering empty arrays, ties, and very unequal lengths.
\solution Swap the arrays if necessary so that $m\le n$. If $m=0$, read the central one or two elements of $B$ directly. Otherwise put
\[
h=\left\lfloor\frac{m+n+1}{2}\right\rfloor.
\]
Take $i$ elements from the beginning of $A$ and $j=h-i$ from the beginning of $B$ as the left half. The feasible search range is
\[
\max(0,h-n)\le i\le\min(m,h).
\]
The four boundary elements are
\[
\begin{array}{ll}
L_A=A_{i-1}\ (i>0),&R_A=A_i\ (i<m),\\
L_B=B_{j-1}\ (j>0),&R_B=B_j\ (j<n).
\end{array}
\]
Use $-\infty$ for a missing left boundary and $+\infty$ for a missing right boundary; these sentinels are used only in comparisons.

A partition is correct exactly when $L_A\le R_B$ and $L_B\le R_A$. Internal sortedness then guarantees that every left-half element is at most every right-half element. If $L_A>R_B$, too many elements came from $A$, so decrease $i$. Larger $i$ cannot work, since $L_A$ is nondecreasing and $R_B$ nonincreasing as $i$ grows. If $L_B>R_A$, increase $i$ by the symmetric argument. Binary search therefore applies.

A correct partition exists: conceptually sort all elements, using any fixed stable order for ties. The first $h$ elements form a prefix of each original array. Every elimination preserves at least one correct partition, so the search terminates. For odd $m+n$, return $\max(L_A,L_B)$. For even $m+n$, return
\[
\frac{\max(L_A,L_B)+\min(R_A,R_B)}2.
\]
In the even, nonempty case both halves contain elements, so the final average uses finite data values. Ties may permit several partitions but all yield the same median. There are at most $m+1$ candidate integers, proving the time and space bounds.
\insight Search for the number of elements contributed by the shorter array to the left half. Monotonicity of the boundary tests replaces the need to merge the data.

\problem{CA03}{Inversions and Minimum Adjacent Swaps}{}
A real array $a_1,\ldots,a_n$ may contain ties. An inversion is a pair $i<j$ with $a_i>a_j$. Count inversions in $O(n\log n)$ time and prove that their number equals the minimum number of adjacent swaps needed to make the array nondecreasing. Evaluate $(3,1,2,1)$ and explain how ties are handled.
\solution Augment merge sort with an inversion count. Recursively sort each half and count its internal inversions. During merging, output the left element whenever it is at most the right element. Otherwise output the right element and add the number of unconsumed left elements: all of them are strictly larger. Each cross-half inversion is counted exactly when its right element is output first, so none is missed or duplicated.

Each merge level costs $O(n)$ and there are $O(\log n)$ levels. A merge buffer uses $O(n)$ additional space. Empty and one-element arrays have count zero. Taking the left element first on equality avoids counting ties as inversions. The counter must accommodate values up to $n(n-1)/2$.

For optimality of the swap count, an adjacent swap leaves the total contribution of every other element unchanged: that element lies on the same side of both swapped positions. Only the swapped pair can change its inversion status, so one swap decreases the total count by at most one. Any sorting sequence therefore needs at least the initial inversion count in swaps.

Conversely, any unsorted array contains a strictly decreasing adjacent pair; otherwise all adjacent comparisons would already imply sortedness. Swapping such a pair decreases the inversion count by exactly one. Repeating sorts the array in exactly the initial inversion count, attaining the lower bound. For the example, the inverted index pairs are
\[
(1,2),(1,3),(1,4),(3,4),
\]
so both the inversion count and the minimum adjacent-swap count are $4$.
\insight Divide and conquer assigns each pair a unique counting step. The same count is a potential function that certifies the optimal number of swaps.
